% Part: computability
% Chapter: recursive-functions
% Section: primitive-recursion

\documentclass[../../../include/open-logic-section]{subfiles}

\begin{document}

\olfileid{cmp}{rec}{pre}
\olsection{العودية البدائية}

كل عدد طبيعي يُنال بالبدء من~$0$
وتكرار عملية اللاحق~$+1$ مرات متناهية؛ فهو إما~$0$
وإما لاحق … لاحق~$0$. وهذه الخاصية تتيح لنا تعيين دالة
$h\colon\Nat \to \Nat$ على وجه مخصوص:
(أ) نضع قيمة $h$ عند الحجة~$0$؛ و(ب) نبيّن كيف نستخرج، من
القيمة $h(x)$، القيمة $h(x+1)$. أما (أ) فيعطينا $h(0)$ من غير واسطة،
فتتعين $h$ عند~$0$. وأما (ب) فنستعملها أولًا عند $x=0$
لنحسب $h(1)=h(0+1)$ من~$h(0)$. ثم نعيد استعمالها عند $x=1$
فنحسب $h(2)=h(1+1)$ من~$h(1)$، وهكذا. ومهما أخذنا عددًا طبيعيًا~$x$، بلغنا بعد خطوات متناهية
موضع تعريف $h(x+1)$ من~$h(x)$، فتكون $h$ معرفة
عند كل $x\in\Nat$.

ومن أمثلة ذلك أن نعيّن $h\colon\Nat\to\Nat$ بالمعادلتين الآتيتين:
\begin{align*}
h(0) & =  1\\
h(x+1) & =  2 \cdot h(x)
\end{align*}
فإن كانت عملية الضرب معلومة لنا، وفت هاتان المعادلتان
بالمطلوب في (أ) و(ب). وإذا أجرينا المعادلة الثانية مرة بعد مرة حصل
\begin{align*}
  h(1) & = 2\cdot h(0) = 2,\\
  h(2) & = 2\cdot h(1) = 2\cdot 2,\\
  h(3) & = 2 \cdot h(2) = 2\cdot 2 \cdot 2,\\
  & \vdots
\end{align*}
فيتبين أن الدالة~$h$ المعينة بهذا الطريق هي $h(x)=2^x$.

ومن الخاصية المتقدمة للأعداد الطبيعية يلزم وجود دالة واحدة لا غير~$h$ تفي
بالشرطين. ونسمّي هاتين المعادلتين \emph{تعريفًا بالعودية
البدائية} للدالة~$h$. أما «العودية» فلأن تعريف~$h$، وبخاصة معادلته
الثانية، يرجع إلى $h$ نفسها في الطرف الأيمن. وأما «البدائية» فلأن
تعريف $h(x+1)$ لا يستعمل من قيم الدالة إلا~$h(x)$، وهي السابقة لها مباشرة.
فهذا أبسط وجوه التعريف العودي على~$\Nat$.

ولا يقتصر هذا الوجه على المثال المتقدم؛ بل نعرّف به ما هو أسبق، كالجمع والضرب.
وهاتان دالتان ثنائيتا الحجة، ولذلك نثبت أحد موضعي
الحجج ونجري العودية على الآخر. فإذا أردنا تعريف $\Add(x,y)$، مثلًا، يمكننا تثبيت
~$x$ وتعريف القيمة أولًا عند $y=0$، ثم تعريفها عند $y+1$ بدلالة~$y$.
وحيث أبقينا $x$ ثابتًا، فإنه يرد في طرفي كل من معادلتي التعريف.
\begin{align*}
\Add(x,0) & =  x\\
\Add(x,y+1) & =  \Add(x,y)+1
\end{align*}
وبهاتين المعادلتين تتعين $\Add$ مهما كانت قيمتا $x$ و~$y$. فلطلب
$\Add(2,3)$، مثلًا، نطبق معادلتي التعريف عند $x=2$: نستعمل الأولى
لإيجاد $\Add(2,0)=2$، ثم نستعمل الثانية تباعًا فنجد
$\Add(2,1)=2+1=3$ و$\Add(2,2)=3+1=4$ و$\Add(2,3)=4+1=5$.

وقد ورد في تعريف $\Add$ الرمز $+$ في الطرف الأيمن من المعادلة
الثانية؛ غير أن استعماله مقصور على زيادة~$1$. فلم نستعمل في الحقيقة إلا دالة اللاحق
$\Succ(z)=z+1$ وطبقناها على القيمة السابقة $\Add(x,y)$ لتعريف
$\Add(x,y+1)$. فكأن التعريف العودي يستند إلى
دالة واحدة نجريها على القيمة السابقة. ولا مانع، بل قد تدعو
الحاجة، إلى أن تعتمد تلك الدالة على غير هذه القيمة؛ فنجعل
$x$ و~$y$ من حججها أيضًا. ومثاله:
\begin{align*}
  \Mult(x,0) & =  0 \\
  \Mult(x,y+1) & =  \Add(\Mult(x,y),x)
\end{align*}
فهذا تعريف للدالة $\Mult$ بالعودية البدائية: نجري فيه الدالة $\Add$ على
القيمة السابقة $\Mult(x,y)$ والحجة الأولى~$x$ جميعًا. وبه
تتعين الدالة~$\Mult(x,y)$ لكل حجتين $x$ و~$y$. ومن ذلك أن قيمة
 $\Mult(2,3)$ تتحدد بحساب $\Mult(2,0)$ ثم $\Mult(2,1)$ ثم $\Mult(2,2)$
ثم~$\Mult(2,3)$ تباعًا:
\begin{align*}
  \Mult(2,0) & = 0\\
  \Mult(2,1) & = \Mult(2,0+1) =
  \Add(\Mult(2,0), 2) = \Add(0, 2) = 2\\
  \Mult(2,2) & = \Mult(2,1+1) =
  \Add(\Mult(2,1), 2) = \Add(2, 2) = 4\\
  \Mult(2,3) & = \Mult(2,2+1) =
  \Add(\Mult(2,2), 2) = \Add(4, 2) = 6
\end{align*}

وعلى هذا يكون النمط العام كما يأتي. إذا أردنا تعريف دالة
~$h(x_0,\dots,x_{k-1},y)$ بالعودية البدائية، أوردنا معادلتين: الأولى تعيّن
$h(x_0,\dots,x_{k-1},0)$ من غير إحالة إلى~$h$. والثانية تعيّن
$h(x_0,\dots,x_{k-1},y+1)$ بدلالة $h(x_0,\dots,x_{k-1},y)$ والحجج
الأخرى $x_0$، …،~$x_{k-1}$ و~$y$. ولا نستعمل في المعادلة الثانية من قيم
الدالة~$h$ إلا السابقة مباشرة. فإن جعلنا العمليتين اللتين يؤديهما
الطرفان الأيمنان دالتين، وسميناهما $f$ و~$g$، صار النمط العام
لتعريف دالة جديدة~$h$ بالعودية البدائية كما يأتي:
\begin{align*}
  h(x_0, \dots, x_{k-1}, 0) & = f(x_0, \dots, x_{k-1})\\
  h(x_0, \dots, x_{k-1}, y+1) & =
  g(x_0, \dots, x_{k-1}, y, h(x_0, \dots, x_{k-1}, y))
\end{align*}
ففي مثال $\Add$ نجد $k=1$ و$f(x_0)=x_0$ (دالة الهوية)، و
$g(x_0,y,z)=z+1$ (الدالة الثلاثية الحجج التي تعيد لاحق حجتها الثالثة):
\begin{align*}
  \Add(x_0, 0) & = f(x_0) = x_0\\
  \Add(x_0, y+1) & = g(x_0, y, \Add(x_0, y)) =
  \Succ(\Add(x_0, y))
\end{align*}
وأما مثال $\Mult$ ففيه $f(x_0)=0$ (الدالة الثابتة التي تعيد~$0$
دائمًا) و$g(x_0,y,z)=\Add(z,x_0)$ (الدالة الثلاثية الحجج التي تعيد
مجموع حجتيها الأخيرة والأولى):
\begin{align*}
  \Mult(x_0,0) & =  f(x_0) = 0 \\
  \Mult(x_0,y+1) & =  g(x_0, y, \Mult(x_0,y)) =
  \Add(\Mult(x_0,y), x_0)
\end{align*}

\end{document}
